%% MLMS Studio — compact method note (Elsevier-style preprint, target $\approx$5 pages).
%% Full manuscript: \texttt{latex/articlelsevier/main.tex}.

\documentclass[preprint,12pt]{elsarticle}
\usepackage{mathtools}
\usepackage{float}
\usepackage{amssymb,natbib}
\usepackage{booktabs}

\journal{Knowledge Based Systems}

\begin{document}

\begin{frontmatter}

\title{MS-ML Studio: methodological core for MS preprocessing, cross-validated feature extraction, and MLP classification}

\address[addr1]{Departamento de Bioinform{\'a}tica, Universidade Federal do Paran{\'a}, Brazil}
\address[addr2]{Setor Litoral, Universidade Federal do Paran{\'a}, Brazil}
\address[addr3]{Departamento de Patologia M{\'e}dica, Universidade Federal do Paran{\'a}, Brazil}

\author[addr1]{Calebe Elias Ribeiro Brima}
\author[addr1]{Marilson Reque}
\author[addr3]{Cibelle B. Dallagassa}
\author[addr3]{Cyntia M. T. Fadel-Picheth}
\author[addr2]{Luciano F. Huergo}
\author[addr1]{Roberto Tadeu Raittz\corref{cor1}}
\cortext[cor1]{Corresponding author}
\ead{raittz@ufpr.br}

\begin{abstract}
\textbf{MS-ML Studio} is an open-source MATLAB application that unifies mass-spectrometry (MS) preprocessing, three complementary feature extractors with embedded K-fold cross-validation, and multilayer perceptron (MLP) classification. This note summarizes the \textbf{algorithmic pipeline}---normalization, reshaping, optional Icoshift alignment, fuzzy moving-average peak detection, correlation- and probability-based peak selection, Watch Points, Feature Resonance with genetic-algorithm (GA) optimization, and MLP training---as implemented in the MS-ML Library and GUI. An illustrative intact-cell MALDI-TOF \textit{Aeromonas} case study~\cite{Surek2010,Croxatto2012} is referenced only to anchor partitions and metrics; claims remain proof-of-concept software validation, not clinical benchmarking.
\end{abstract}

\begin{keyword}
mass spectrometry \sep feature extraction \sep genetic algorithm \sep MALDI-TOF \sep machine learning \sep open-source software
\end{keyword}

\end{frontmatter}

\section{Introduction}
MS workflows recur through import, preprocessing, feature extraction, and classification~\cite{Krishnamurthy1996,Croxatto2012,Pluskal2010}. MS-ML Studio targets users who need these stages in one graphical session without scripting. The present text \textbf{focuses on methods}; narrative introduction, extended related work, and discussion appear in the full article under \texttt{latex/articlelsevier/}.

\section{Materials and methods}
\label{sec:methods}

\subsection{Reference data (illustrative)}
Spectra follow the intact-cell MALDI-TOF \textit{Aeromonas} setting described by Surek \textit{et al.}~\cite{Surek2010}. For each one-against-all experiment, 90\% of spectra train feature extractors and the MLP; 10\% are held out for blind testing (Table~\ref{tab:partition}).

\subsection{Architecture}
The MS-ML Library implements all algorithms; the GUI orchestrates preprocessing, peak selection, feature extraction with cross-validation, and classification.

\subsection{Preprocessing}

\subsubsection{Import and normalization}
Spectra are two-column text files (m/z, intensity). Min--max normalization maps intensities to $[0,1]$:
\begin{equation}
N(x, x_{\max}, x_{\min}) = \frac{x_i - x_{\min}}{x_{\max} - x_{\min}}\,.
\label{eq:norm}
\end{equation}

\subsubsection{Reshaping}
Each spectrum is binned onto $W$ evenly spaced grid points over the global m/z range; per-bin intensity is the maximum of contributing points and m/z is the mean of their positions~\cite{Liu2013} (full discrete definition in the long manuscript).

\subsubsection{Alignment}
Optional alignment uses Icoshift~\cite{Tomasi2011} on user-selected m/z windows after reshaping.

\subsubsection{Fuzzy moving average (FMA) peaks}
At position $p$, a triangular-window local average $S(x;p,w)$ (Eq.~\ref{eq:fma}) highlights local maxima as candidate peaks; $w$ is half-width and $T$ the triangular weight (same functional form as Watch Points below).
\begin{equation}
S(x; p, w) = \frac{\sum_{i=p-w}^{p+w} x_i \cdot T(2w+1, i)}{2w+1}\,.
\label{eq:fma}
\end{equation}

\subsection{Best peaks selection}
From FMA candidates, \textbf{correlation-based} selection scores each peak's presence/absence vector against class labels (Pearson); peaks with $|\rho|>\tau$ survive. \textbf{Probability-based} selection uses class-conditional presence probabilities for discriminancy and exclusivity (see full text).

\subsection{Feature extraction with cross-validation}
All GA-driven extractors embed K-fold cross-validation: a configuration is retained only if the mean of training and validation correlations improves~\cite{Liu2013}.

\subsubsection{Watch Points}
For odd window length $L$, triangular weights are
\begin{equation}
w(L,n) = \begin{cases} \dfrac{2n}{L+1} & 1 \leq n \leq \dfrac{L+1}{2} \\[\medskipamount]
2 - \dfrac{2n}{L+1} & \dfrac{L+1}{2}+1 \leq n \leq L \end{cases}\,.
\end{equation}
(Even $L$ uses the paired formulas in the repository manuscript.) Feature $W_k$ for watch point $k$ is the dot product of weights with the binary peak vector; the GA optimizes half-widths $\{L_k\}$.

\subsubsection{Feature Resonance}
Peaks are encoded via a sum of $M$ sinusoids with GA-tuned parameters:
\begin{equation}
S(x; \mathbf{a}, \mathbf{b}, \mathbf{c}) = \sum_{i=1}^{M} a_i \sin(b_i + c_i x)\,.
\label{eq:resonance}
\end{equation}
Fitness is the mean absolute Pearson correlation between $S$ and labels; variants operate on m/z grids, best-peak intensities, m/z values, or watch-point vectors.

\subsection{Classification and metrics}
Features from enabled extractors are concatenated and fed to an MLP (Levenberg--Marquardt). Accuracy, precision, recall, and F1 follow standard definitions~\cite{Powers2011}. ARFF export enables external classifiers (e.g.\ WEKA).

\section{Results (summary)}
\label{sec:results}

\begin{table}[H]
  \centering
  \caption{Reference partitioning (same as full manuscript).}
  \label{tab:partition}
  \begin{tabular}{@{}lccc@{}}
    \toprule
    Species & Training ($n$) & Test ($n$) & Total ($n$) \\
    \midrule
    \textit{A.~hydrophila} & 511 & 56 & 567 \\
    \textit{A.~caviae} & 511 & 56 & 567 \\
    \textit{A.~veronii}/\textit{sobria} & 286 & 30 & 316 \\
    \textit{A.~trota} & 511 & 56 & 567 \\
    \bottomrule
  \end{tabular}
\end{table}

Held-out test accuracies and F1-scores on this reference task are reported in the full article; the smallest test fold (\textit{A.~veronii}/\textit{sobria}, $n=30$) warrants cautious interpretation.

\section{Conclusion}
MS-ML Studio implements a closed pipeline from text spectra through GA-regularized feature extractors to MLP classification. This compact note documents the methodological core; extended results, discussion, limitations, and availability match \texttt{latex/articlelsevier/main.tex}.

\section*{Acknowledgements}
Dataset acknowledgment follows Surek \textit{et al.}~\cite{Surek2010}; institutional support as in the full manuscript.

\bibliographystyle{elsarticle-num}
\bibliography{../articlelsevier/mendeley,../articlelsevier/msml_supp}

\end{document}
